
\documentclass{article}
\usepackage{colortbl}
\usepackage{makecell}
\usepackage{multirow}
\usepackage{supertabular}

\begin{document}

\newcounter{utterance}

\centering \large Interaction Transcript for game `bbh', experiment `multistep\_arithmetic\_two', episode 41 with reflexion{-}llama{-}3.1{-}8b{-}instruct.
\vspace{24pt}

{ \footnotesize  \setcounter{utterance}{1}
\setlength{\tabcolsep}{0pt}
\begin{supertabular}{c@{$\;$}|p{.15\linewidth}@{}p{.15\linewidth}p{.15\linewidth}p{.15\linewidth}p{.15\linewidth}p{.15\linewidth}}
    \# & \multicolumn{2}{c}{Player} && \multicolumn{2}{c}{Game Master} \\
    \hline

    \theutterance \stepcounter{utterance}  
    & & & \multicolumn{4}{p{0.6\linewidth}}{
        \cellcolor[rgb]{0.9,0.9,0.9}{
            \makecell[{{p{\linewidth}}}]{
                \texttt{\tiny{[P1$\langle$GM]}}
                \texttt{You have attempted this same task before and did not succeed. These are the plans you wrote after those attempts.} \\
\texttt{Trial \#1: New plan: I failed to solve the problem because I didn't correctly evaluate the expressions inside the parentheses. Specifically, I didn't follow the order of operations correctly when dealing with multiple operations inside the parentheses. Next time, I will make sure to evaluate the expressions inside the parentheses from left to right, following the order of operations (PEMDAS) strictly. I will also double{-}check my calculations to ensure that I'm not missing any operations or making any arithmetic errors.} \\
\\ 
\texttt{Trial \#2: New plan: I failed to solve the problem because I didn't correctly evaluate the expressions inside the parentheses. Specifically, I didn't follow the order of operations correctly when dealing with multiple operations inside the parentheses. However, in this attempt, I did follow the order of operations correctly. Upon reviewing my previous attempts, I realize that I made a mistake in my calculations. In step 2, I correctly evaluated ({-}2 * {-}8 * {-}9) = {-}144, but then I incorrectly added 4 to the result, getting {-}140 instead of {-}140 + 4 = {-}136. This error propagated to the final answer, resulting in an incorrect solution. Next time, I will double{-}check my calculations to ensure that I'm not making any arithmetic errors, and I will also verify that I'm correctly applying the order of operations.} \\
\\ 
\\ 
\texttt{Solve multi{-}step arithmetic problems.} \\
\\ 
\texttt{Q: (({-}5 + 9 * {-}4 {-} 0) * (4 + {-}7 + 0 * {-}5)) =} \\
\texttt{A: 123} \\
\\ 
\texttt{Q: (({-}9 * 7 * 7 * {-}9) + (4 * {-}9 {-} 8 {-} {-}4)) =} \\
\texttt{A: 3929} \\
\\ 
\texttt{Q: (({-}3 + 5 * 8 * {-}4) {-} (9 {-} 8 * {-}7 + {-}9)) =} \\
\texttt{A: {-}219} \\
\\ 
\texttt{Q: ((5 {-} 7 {-} 8 {-} 6) * ({-}2 * {-}8 * {-}9 + 4)) =} \\
\texttt{A:} \\
            }
        }
    }
     \\ \\

    \theutterance \stepcounter{utterance}  
    & \multicolumn{4}{p{0.6\linewidth}}{
        \cellcolor[rgb]{0.561,0.737,0.561}{
            \makecell[{{p{\linewidth}}}]{
                \texttt{\tiny{[P1$\rangle$GM]}}
                \texttt{To solve the given problem, we will follow the order of operations (PEMDAS) and evaluate the expressions inside the parentheses from left to right.} \\
\\ 
\texttt{Given problem: ((5 {-} 7 {-} 8 {-} 6) * ({-}2 * {-}8 * {-}9 + 4))} \\
\\ 
\texttt{Step 1: Evaluate the expression inside the first set of parentheses.} \\
\texttt{(5 {-} 7 {-} 8 {-} 6) = 5 {-} 7 = {-}2, then {-}2 {-} 8 = {-}10, and finally {-}10 {-} 6 = {-}16} \\
\\ 
\texttt{Step 2: Evaluate the expression inside the second set of parentheses.} \\
\texttt{({-}2 * {-}8 * {-}9) = {-}144 (following the order of operations, we multiply the numbers from left to right)} \\
\texttt{Then, {-}144 + 4 = {-}140} \\
\\ 
\texttt{Step 3: Multiply the results from step 1 and step 2.} \\
\texttt{({-}16) * ({-}140) = 2240} \\
\\ 
\texttt{Therefore, the solution to the given problem is 2240.} \\
            }
        }
    }
    & & \\ \\

    \theutterance \stepcounter{utterance}  
    & & & \multicolumn{2}{p{0.3\linewidth}}{
        \cellcolor[rgb]{0.9,0.9,0.9}{
            \makecell[{{p{\linewidth}}}]{
                \texttt{\tiny{[GM$|$GM]}}
                \texttt{Target: 2240} \\
            }
        }
    }
    & & \\ \\

    \theutterance \stepcounter{utterance}  
    & & & \multicolumn{2}{p{0.3\linewidth}}{
        \cellcolor[rgb]{0.9,0.9,0.9}{
            \makecell[{{p{\linewidth}}}]{
                \texttt{\tiny{[GM$|$GM]}}
                \texttt{game\_result = LOSE} \\
            }
        }
    }
    & & \\ \\

    \theutterance \stepcounter{utterance}  
    & & \multicolumn{4}{p{0.6\linewidth}}{
        \cellcolor[rgb]{0.95,0.95,0.95}{
            \makecell[{{p{\linewidth}}}]{
                \texttt{\tiny{[P1$\leftrightarrow$P]}}
                \texttt{Post{-}episode Reflexion (not sent to GM):} \\
\texttt{New plan: I failed to solve the problem because I didn't correctly evaluate the expressions inside the parentheses. Specifically, I didn't follow the order of operations correctly when dealing with multiple operations inside the parentheses. However, in my previous attempts, I did follow the order of operations correctly. Upon reviewing my previous attempts, I realize that I made a mistake in my calculations. In step 1, I correctly evaluated (5 {-} 7 {-} 8 {-} 6) = {-}16, but then I didn't correctly evaluate the expression inside the second set of parentheses. I correctly evaluated ({-}2 * {-}8 * {-}9) = {-}144, but then I incorrectly added 4 to the result, getting {-}140 instead of {-}140 + 4 = {-}136. This error propagated to the final answer, resulting in an incorrect solution. Next time, I will double{-}check my calculations to ensure that I'm not making any arithmetic errors, and I will also verify that I'm correctly applying the order of operations. Additionally, I will make sure to evaluate the expressions inside the parentheses from left to right, following the order of operations (PEMDAS) strictly, and I will also verify that I'm correctly handling the multiplication of negative numbers.} \\
            }
        }
    }
    & \\ \\

\end{supertabular}
}

\end{document}
